\documentclass[aps,prb,onecolumn,nofootinbib,superscriptaddress]{revtex4-2}
\usepackage{amsmath,amssymb,amsthm,mathtools,bm}
\usepackage{booktabs}
\usepackage[colorlinks=true,linkcolor=blue,citecolor=blue,urlcolor=blue]{hyperref}
\usepackage{microtype}
\newcommand{\One}{\hat{\mathbf{1}}}
\newcommand{\Id}{\mathbf{1}}
\newcommand{\Tr}{\operatorname{Tr}}
\newcommand{\tr}{\operatorname{tr}}
\newcommand{\ket}[1]{\lvert #1\rangle}
\newcommand{\bra}[1]{\langle #1\rvert}
\newcommand{\eps}{\varepsilon}
\newcommand{\hc}{\mathrm{h.c.}}
\hypersetup{pdftitle={Gain and Loss as Singular Limits of Gaussian Transfer Matrices},
pdfauthor={Asadullah Bhuiyan}}
\begin{document}
\title{Gain and Loss as Singular Limits of Gaussian Transfer Matrices}
\author{Asadullah Bhuiyan}
\noaffiliation
\date{October 7, 2026}
\begin{abstract}
Fermion creation and annihilation preserve Gaussian conditional states but
are nilpotent, so they have no invertible conjugation action on fermion
operators. We resolve this apparent mismatch by representing each jump as
a projectively rescaled limit of an invertible odd Gaussian operator. Its
polar decomposition is a weak occupation measurement followed by a
Majorana flip; its action on Majoranas lies in the disconnected component
of $O(2N,\mathbb C)$. Outcome-dependent feedback gives finite-strength
instruments whose limits are precisely the fill and deplete steps of a
fresh-ancilla adaptive circuit. Ordered products have a well-defined
normalized state limit for every fixed finite, nonzero branch, even when
their transfer matrices diverge. We derive the exact Slater-projector
updates at the endpoint and explain why regulator singularities, physical
trajectory singular values, and tangent Lyapunov rates must be kept distinct.
\end{abstract}
\maketitle
\setcounter{tocdepth}{2}
\tableofcontents

\section{The physical question and conventions}
An adaptive occupation measurement may have to add or remove a fermion
after detecting the wrong occupation. Such an operation is Gaussian in
its action on states, but cannot be written as an invertible exponential
of a quadratic operator: a creation or annihilation operator squares to
zero. The useful statement is not that the jump itself belongs to a
classical matrix group. Rather, it is a singular limit of operators whose
Majorana conjugation actions belong to that group. The scalar used to
take the limit and the state on which the operator acts are indispensable.

We use $N$ complex modes $\hat\psi_i$ with
$\{\hat\psi_i,\hat\psi_j^\dagger\}=\delta_{ij}$ and
$\{\hat\psi_i,\hat\psi_j\}=0$. A normalized column $\chi\in\mathbb C^N$
defines
\begin{equation}
 \hat\chi^\dagger=\sum_i\chi_i\hat\psi_i^\dagger,\qquad
 \hat\chi=\sum_i\chi_i^*\hat\psi_i,\qquad
 \hat n=\hat\chi^\dagger\hat\chi,\qquad \hat n^2=\hat n.
 \label{eq:mode}
\end{equation}
Many-body operators carry hats. Following Bhuiyan, Pan, and Jian
(BPJ)~\cite{BhuiyanPanJian2026}, the correlation matrix is
$G_{ij}=\Tr(\hat\rho\hat\psi_i^\dagger\hat\psi_j)$. For occupied columns
in the convention~\eqref{eq:mode}, it is convenient to use
\begin{equation}
 C=G^{\mathsf T},\qquad C_{ij}=\Tr(\hat\rho\hat\psi_j^\dagger\hat\psi_i).
 \label{eq:C-convention}
\end{equation}
Thus a Slater determinant with orthonormal occupied columns $F$ has
$C=FF^\dagger$; no centered covariance is denoted by $G$ here.

The Majoranas
\begin{equation}
 \hat\gamma_{2i-1}=\hat\psi_i+\hat\psi_i^\dagger,\qquad
 \hat\gamma_{2i}=-i(\hat\psi_i-\hat\psi_i^\dagger),\qquad
 \{\hat\gamma_a,\hat\gamma_b\}=2\delta_{ab}
 \label{eq:majoranas}
\end{equation}
span a $2N$-dimensional space. We allow both even Gaussian operators and
their products with Majorana operators, which are odd. This extends the
invertible quadratic-exponential description by its odd Clifford sector;
the state and measurement formalism is consistent with fermionic linear
optics~\cite{Bravyi2005}. A physical implementation of an odd system branch
exchanges parity with an ancillary fermion, rather than violating total
fermion-parity conservation.

\section{From a nilpotent jump to an orthogonal transformation}
\subsection{An invertible regulator and its polar decomposition}
For $0<\eps\leq1$, define
\begin{align}
 \hat j_+(\eps)&=\hat\chi^\dagger+\eps\hat\chi,&
 \hat j_-(\eps)&=\hat\chi+\eps\hat\chi^\dagger,\nonumber\\
 \hat v_\pm(\eps)&=\eps^{-1/2}\hat j_\pm(\eps).
 \label{eq:regulators}
\end{align}
The canonical anticommutation relations give
\begin{equation}
 \hat j_\pm^2=\eps\One,\qquad \hat v_\pm^2=\One,\qquad
 \sqrt\eps\,\hat v_+\longrightarrow\hat\chi^\dagger,\quad
 \sqrt\eps\,\hat v_-\longrightarrow\hat\chi.
 \label{eq:limits}
\end{equation}
Every nonzero regulator therefore supplies an inverse. Exact gain and
loss do not: $\hat\chi^2=(\hat\chi^\dagger)^2=0$.

To identify the physical content, set
\begin{equation}
 \hat\Gamma=\hat\chi+\hat\chi^\dagger,\qquad
 \hat\Gamma^\dagger=\hat\Gamma,\qquad \hat\Gamma^2=\One,\qquad
 \lambda=-\tfrac12\log\eps.
\end{equation}
Direct multiplication gives
\begin{align}
 \hat v_+^\dagger\hat v_+
 &=\eps^{-1}(\One-\hat n)+\eps\hat n
   =e^{-2\lambda(2\hat n-\One)},\nonumber\\
 \hat v_-^\dagger\hat v_-
 &=\eps(\One-\hat n)+\eps^{-1}\hat n
   =e^{+2\lambda(2\hat n-\One)}.
 \label{eq:positive-squares}
\end{align}
Since $\hat\Gamma\hat n=\hat\chi$ and
$\hat\Gamma(\One-\hat n)=\hat\chi^\dagger$, their polar decompositions are
\begin{equation}
 \hat v_\pm=\hat\Gamma\,e^{\mp\lambda(2\hat n-\One)}.
 \label{eq:polar}
\end{equation}
The positive factor filters occupation; the odd unitary factor flips it.
The filter has the form of the weak single-mode measurement used in
BPJ Appendix B.2~\cite{BhuiyanPanJian2026}. A scalar normalization and a
complete set of outcomes are still required to make a physical instrument.

\subsection{The group and the meaning of its boundary}
For any complex column $z$, let
$\hat\mu(z)=\sum_a z_a\hat\gamma_a$. The relevant quadratic form is
bilinear, not Hermitian:
\begin{equation}
 \hat\mu(z)^2=(z^{\mathsf T}z)\One.
\end{equation}
When $z^{\mathsf T}z\ne0$, anticommuting $\hat\gamma_a$ through
$\hat\mu(z)$ gives
\begin{equation}
 \hat\mu(z)\hat\gamma_a\hat\mu(z)^{-1}
 =\sum_b t[z]_{ab}\hat\gamma_b,\qquad
 t[z]=-\Id+\frac{2zz^{\mathsf T}}{z^{\mathsf T}z}.
 \label{eq:reflection}
\end{equation}
The rank-one factor in this expression is an idempotent, hence
\begin{equation}
 t[z]^{\mathsf T}t[z]=\Id,\qquad t[z]^2=\Id,\qquad
 \det t[z]=(-1)^{2N-1}=-1.
 \label{eq:orthogonal}
\end{equation}
Thus an invertible odd Clifford operator induces a single-particle transfer
matrix (sTM) in the disconnected component of $O(2N,\mathbb C)$. On the
two Majoranas of the measured mode,
\begin{equation}
 z_+(\eps)=\tfrac12
 \begin{pmatrix}\eps^{-1/2}+\eps^{1/2}\\
 i(\eps^{1/2}-\eps^{-1/2})\end{pmatrix},\qquad
 z_+^{\mathsf T}z_+=1,
 \label{eq:regulated-z}
\end{equation}
represents $\hat v_+$. The creation operator instead corresponds to
$z_0=(1,-i)^{\mathsf T}/2$, with $z_0^{\mathsf T}z_0=0$. The jump is a
null Clifford direction, so Eq.~\eqref{eq:reflection} cannot be evaluated
there by ordinary matrix substitution.

The same degeneration is transparent in the local Nambu basis:
\begin{equation}
 \hat v_+\begin{pmatrix}\hat\chi\\\hat\chi^\dagger\end{pmatrix}\hat v_+^{-1}
 =\begin{pmatrix}0&\eps^{-1}\\\eps&0\end{pmatrix}
 \begin{pmatrix}\hat\chi\\\hat\chi^\dagger\end{pmatrix}.
 \label{eq:nambu}
\end{equation}
Orthogonal spectator modes acquire a minus sign under this odd
conjugation. In an orthonormal mode basis the full sTM has singular values
$\{\eps^{-1},1,\ldots,1,\eps\}$, although all its eigenvalues are $\pm1$.
The divergent singular values are not divergent eigenvalues.

Here ``boundary'' means a boundary of the projective operator or
normalized-state description: $[\hat v_+]\to[\hat\chi^\dagger]$, with
the rescaling specified in Eq.~\eqref{eq:limits}. It is \emph{not} a finite
singular matrix in the ordinary closure of $O(2N,\mathbb C)$. A finite
limit of matrices satisfying $t^{\mathsf T}t=\Id$ still satisfies that
identity and is invertible. In the present limit the conjugation matrix
escapes to infinity, while its rescaled many-body representative has a
finite nonzero limit. Normalized Choi states supply a concrete way to
retain that limit, as developed below.

\section{A physical measurement and the adaptive feedback choices}
Define two positive weak-measurement operators,
\begin{equation}
 \hat M_0=\frac{\One-\hat n+\eps\hat n}{\sqrt{1+\eps^2}},\qquad
 \hat M_1=\frac{\hat n+\eps(\One-\hat n)}{\sqrt{1+\eps^2}}.
 \label{eq:weak-measurement}
\end{equation}
The subscripts label outcomes that become ``empty'' and ``occupied'' in
the strong limit. Their effects obey
\begin{equation}
 \hat E_0=\frac{\One-\hat n+\eps^2\hat n}{1+\eps^2},\qquad
 \hat E_1=\frac{\hat n+\eps^2(\One-\hat n)}{1+\eps^2},\qquad
 \hat E_0+\hat E_1=\One.
 \label{eq:effects}
\end{equation}
This is a positive-operator-valued measure (POVM). In the weighted
convention $\sum_m w_m\hat K_m^\dagger\hat K_m=\One$, our displayed
Kraus operators absorb $\sqrt{w_m}$, so all remaining weights equal one.
For example,
$p_0=[1-\langle\hat n\rangle+\eps^2\langle\hat n\rangle]/(1+\eps^2)$.

Outcome-dependent unitary feedback leaves these effects unchanged. The
following three choices have distinct physical purposes:
\begin{center}
\begin{tabular}{llll}
\toprule
Feedback & $\hat K_0$ & $\hat K_1$ & $\eps\to0$ endpoint\\
\midrule
Fill & $\hat\Gamma\hat M_0$ & $\hat M_1$
 & $\{\hat\chi^\dagger,\hat n\}$\\
Deplete & $\hat M_0$ & $\hat\Gamma\hat M_1$
 & $\{\One-\hat n,\hat\chi\}$\\
Flip both & $\hat\Gamma\hat M_0$ & $\hat\Gamma\hat M_1$
 & $\{\hat\chi^\dagger,\hat\chi\}$\\
\bottomrule
\end{tabular}
\end{center}
In particular,
\begin{equation}
 \hat\Gamma\hat M_0=\frac{\hat j_+(\eps)}{\sqrt{1+\eps^2}},\qquad
 \hat\Gamma\hat M_1=\frac{\hat j_-(\eps)}{\sqrt{1+\eps^2}}.
 \label{eq:normalized-jumps}
\end{equation}
Every listed finite-$\eps$ branch is invertible and Gaussian. The first
two instruments approach exactly the target-fill and target-empty
measurement-and-correction steps, up to irrelevant branch phases. They
are the fresh-ancilla idealization of the occupation-measurement and
fermionic-swap (fSWAP) construction
in BPJ Sec.~III~\cite{BhuiyanPanJian2026}, not its finite-bath dynamics.
The third instrument flips either result and must not be substituted for
targeted feedback. Pairing gain with loss is therefore one completion,
not a unique requirement. At $\eps=1$ no occupation information is gained;
these instruments reduce to random unitary evolution.

Equivalently, the flip-both pair uses the unnormalized $\hat v_\pm$ with
equal weights
\begin{equation}
 w_+=w_-=\frac1{\eps+\eps^{-1}}=\frac1{2\cosh(2\lambda)}.
\end{equation}
This recovers the scalar that conjugation matrices discard. Multiplying
operators along a selected record is legitimate, but does not specify a
Born-sampled ensemble until the instrument and its branch probabilities
are fixed.

There are two symmetry distinctions. First, finite-$\eps$ jump branches
superpose charge raising and lowering and generally break $U(1)$
covariance. At $\eps=0$ every branch has a definite charge change, so
$e^{i\theta\hat N_F}\hat\chi^{\dagger}e^{-i\theta\hat N_F}
=e^{i\theta}\hat\chi^\dagger$ and its conjugate make the branch map
$\hat\rho\mapsto\hat K\hat\rho\hat K^\dagger$ $U(1)$ covariant.
The regulator is an auxiliary deformation, not a claim that the original
protocol leaves its symmetry class. Equality of the measurement effects
with a BPJ construction alone does not classify all feedback-dressed
many-body evolution operators (mEOs).

Second, an odd Kraus operator does not create a parity-odd density
operator. If $\hat P=(-1)^{\hat N_F}$, $\hat P\hat K\hat P=-\hat K$, and
$\hat P\hat\rho\hat P=\hat\rho$, then
\begin{equation}
 \hat P(\hat K\hat\rho\hat K^\dagger)\hat P
 =\hat K\hat\rho\hat K^\dagger.
 \label{eq:parity}
\end{equation}
A definite-parity state can switch sectors without acquiring coherence
between them. A globally parity-preserving dilation supplies the exchanged
parity through an ancilla. No extension of symmetry-class POVM no-go
theorems is needed for these elementary statements.

\section{Ordered products and a controlled singular limit}
\subsection{Composition before the limit}
Fix a finite record $\boldsymbol m$ and write its chronological operator
as
\begin{equation}
 \hat V_{\boldsymbol m}=\hat O_q\cdots\hat O_1,
 \label{eq:word}
\end{equation}
where the rightmost factor acts first. The factors can include gain,
loss, occupation projectors, and invertible charge-conserving Gaussian
operations. Any feedback schedule is held fixed by the record.
Replace each jump by $\hat j_\pm(\eps_k)$ and each projector by
$\One-\hat n+\eps_k\hat n$ or
$\hat n+\eps_k(\One-\hat n)$ as appropriate. At every nonzero regulator
the product $\hat W(\bm\eps)$ is invertible and Gaussian.

With BPJ's row-index convention,
\begin{equation}
 \hat V\hat\gamma_a\hat V^{-1}=\sum_b t[\hat V]_{ab}\hat\gamma_b,
 \qquad
 t[\hat V_2\hat V_1]=t[\hat V_1]t[\hat V_2].
 \label{eq:order}
\end{equation}
The order reversal follows by applying $\hat V_2$ to each operator in
$\hat V_1\hat\gamma_a\hat V_1^{-1}$. Thus $t[\hat W]=t[\hat O_1]\cdots
t[\hat O_q]$ with the regulated factors understood; transposed sTMs
instead multiply in the same order as Eq.~\eqref{eq:word}. Scalar branch
weights do not appear in either convention. An even quadratic factor has
determinant $+1$; a product with $r$ odd factors has determinant $(-1)^r$.

The product must be formed with its chronological data intact. Already
on one mode,
\begin{equation}
 \hat\chi\hat\chi^\dagger=\One-\hat n,\qquad
 \hat\chi^\dagger\hat\chi=\hat n,
 \label{eq:order-example}
\end{equation}
whereas opposite diagonal ``zero/pole'' factors can multiply to the
identity and erase this distinction. At finite regulator the full Nambu
matrix, including its anomalous blocks, is needed. If its blocks are
$\bigl(\begin{smallmatrix}A&B\\C&D\end{smallmatrix}\bigr)$, the upper-left
block of a product contains $A_1A_2+B_1C_2$, not just $A_1A_2$.

\subsection{What converges, and under which assumptions}
For a fixed finite number of factors and fixed mode wavefunctions,
polynomial continuity gives, in any finite-dimensional operator norm,
\begin{equation}
 \hat W(\bm\eps)=\hat V_{\boldsymbol m}+O(\eta),\qquad
 \eta=\max_k|\eps_k|.
 \label{eq:word-convergence}
\end{equation}
The same statement holds with the physical normalizations in
Eqs.~\eqref{eq:weak-measurement} and \eqref{eq:normalized-jumps}.
For an input with $\|\hat V_{\boldsymbol m}\ket{\psi}\|>0$, normalization
is continuous, so the normalized branch converges independently of
regulator ratios. The constants depend on the word and its nonzero
success probability; this is not a uniform statement at large depth or
near a forbidden outcome.

For an operator-level version choose a normalized, maximally entangled
Gaussian reference with $N$ reference modes and $N$ physical modes,
\begin{equation}
 \ket{\mathrm{REF}}=2^{-N/2}\prod_{i=1}^N
 (\hat\psi_{i,R}^\dagger+\hat\psi_{i,P}^\dagger)\ket{0},\qquad
 \|\hat V_P\ket{\mathrm{REF}}\|^2=2^{-N}\Tr(\hat V^\dagger\hat V).
 \label{eq:choi-reference}
\end{equation}
A fixed fermionic mode ordering is implicit. If $\hat V\ne0$, its
normalized Choi vector and full covariance are therefore limits of those
of $\hat W$, irrespective of how the regulators approach zero. This is
a concrete state-level meaning of the boundary, with no assertion that
a chosen homogeneous transfer chart remains finite or invertible.

There is a useful sharper rate for \emph{jump-only regularization}, with
charge-conserving invertible factors between jumps and a fixed-charge
input (including the reference above). Flipping one regulated jump from
gain to loss or vice versa changes the output charge by two. Thus
\begin{equation}
 \ket{\Psi(\bm\eps)}=\ket{\Psi_0}
       +\sum_k\eps_k\ket{\Psi_k}+O(\eta^2)
\end{equation}
has no first-order cross term in its norm or in any charge-neutral
bilinear expectation. For a nonzero limiting branch,
\begin{equation}
 G(\bm\eps)-G(0)=O(\eta^2),\qquad
 \langle\hat\psi_i\hat\psi_j\rangle_{\bm\eps}=O(\eta).
 \label{eq:rates}
\end{equation}
The \emph{normal} correlation block converges quadratically; the full
Majorana covariance generally converges only linearly. If projectors are
also regularized, their first-order corrections need not change charge,
so only the general $O(\eta)$ bound is asserted.

The nonzero-branch hypothesis cannot be removed. Two same-mode gains
give zero exactly, but their regulated product is
\begin{equation}
 (\hat\chi^\dagger+\eps_2\hat\chi)
 (\hat\chi^\dagger+\eps_1\hat\chi)
 =\eps_1\hat n+\eps_2(\One-\hat n).
 \label{eq:zero-branch}
\end{equation}
Normalizing this vanishing operator can retain the arbitrary ratio
$\eps_1/\eps_2$. There is no exact physical branch to which that normalized
state should converge.

\section{The exact endpoint without a regulator}
The regularization explains the geometry, but exact state propagation
need not use ill-conditioned matrices. Let a normalized pure Slater
state have occupied projector $C=FF^\dagger$, with $F^\dagger F=\Id$.
For a gain, decompose its column as
$\chi=C\chi+(\Id-C)\chi$. Creation in the occupied component vanishes by
Pauli exclusion. The surviving vector is orthogonal to the old occupied
space, so it adds one normalized column. A loss removes the occupied
component $C\chi$. Consequently,
\begin{align}
 p_+&=\chi^\dagger(\Id-C)\chi,&
 C_+&=C+\frac{(\Id-C)\chi\chi^\dagger(\Id-C)}{p_+},\nonumber\\
 p_-&=\chi^\dagger C\chi,&
 C_-&=C-\frac{C\chi\chi^\dagger C}{p_-}.
 \label{eq:exact-updates}
\end{align}
These expressions require $p_\pm>0$; a zero denominator signals a
vanishing branch, not a prescription for a pseudoinverse. They give
$C_\pm^2=C_\pm$, $\tr C_\pm=\tr C\pm1$,
$C_+\chi=\chi$, and $C_-\chi=0$. They are the pure-state specialization
of Gaussian measurement propagation~\cite{Bravyi2005}.

The two projective branches needed alongside the jumps follow from
$\hat n=\hat\chi^\dagger\hat\chi$ and
$\One-\hat n=\hat\chi\hat\chi^\dagger$. After loss the mode is empty,
so refilling it adds $P_\chi=\chi\chi^\dagger$; after gain the mode is
filled, so emptying it subtracts the same projector. Therefore
\begin{equation}
 C_{\hat n}=C_-+P_\chi,\quad p_{\hat n}=p_-,\qquad
 C_{\One-\hat n}=C_+-P_\chi,\quad p_{\One-\hat n}=p_+.
 \label{eq:projective-updates}
\end{equation}
Equations~\eqref{eq:exact-updates} and \eqref{eq:projective-updates}
handle all four outcomes of the exact fill/deplete instruments using only
one-particle matrices. The second step of either projector factorization
has conditional probability one, explaining the displayed probabilities.

For an intervening even operator
$\hat V=\exp[-\sum_{ij}M_{ij}\hat\psi_i^\dagger\hat\psi_j]$, define
$A=e^{-M}=t^{-1}$. It sends the occupied columns to $AF$, giving
\begin{equation}
 C'=AF(F^\dagger A^\dagger AF)^{-1}F^\dagger A^\dagger,
 \qquad
 \|\hat V\ket F\|^2=\det(F^\dagger A^\dagger AF).
 \label{eq:even-update}
\end{equation}
An additional scalar multiplying $\hat V$ contributes its squared modulus
to the norm. These formulas follow from the Gram matrix of the transformed
Slater columns, and do not assume $A$ is unitary. Using
Eq.~\eqref{eq:exact-updates} after each even step automatically retains
chronology and Pauli blocking. The product of successive conditional norm
factors is the norm squared of the complete word; those factors are Born
probabilities only when the operations belong to the specified instrument.

\section{Consequences for trajectory spectra}
The group description does not by itself identify a physical Lyapunov
spectrum. A single regulated jump has logarithmic sTM singular values
$\pm\log(1/\eps)$, but these are regulator scales, not long-time rates.
For a particularly simple counterexample, repeating the same normalized
odd element twice gives $\hat v_+^2=\One$ and $t[\hat v_+]^2=\Id$.
At fixed $\eps>0$ an arbitrarily long repetition has zero asymptotic sTM
growth rates even though one step becomes arbitrarily ill-conditioned
as $\eps\to0$. For the rescaled operator,
$\hat j_+^2=\eps\One$, so the many-body scalar has a different fate.
The regulator limit and the long-record limit cannot be interchanged
on the strength of Eq.~\eqref{eq:word-convergence}.

The physical endpoint can instead be characterized directly. For a
nonzero exact Gaussian record with a definite charge shift, propagate a
maximally mixed input and set
\begin{equation}
 \hat\rho_{\boldsymbol m}=\frac{\hat V_{\boldsymbol m}
 \hat V_{\boldsymbol m}^\dagger}{Z_{\boldsymbol m}},\qquad
 Z_{\boldsymbol m}=\Tr(\hat V_{\boldsymbol m}\hat V_{\boldsymbol m}^\dagger)
                 =2^N p_{\boldsymbol m}.
 \label{eq:maxmix}
\end{equation}
The endpoint is number-conserving Gaussian. If its occupation eigenvalues
are $\nu_a$, natural-orbital factorization gives its Fock-space eigenvalues
and hence the squared singular values of the mEO:
\begin{equation}
 \sigma_{\bm n}(\hat V_{\boldsymbol m})^2
 =Z_{\boldsymbol m}\prod_a\nu_a^{n_a}(1-\nu_a)^{1-n_a},
 \qquad n_a\in\{0,1\}.
 \label{eq:fock-spectrum}
\end{equation}
The exact occupations $0$ and $1$ encode genuine zero singular values;
the preferred occupation factor equals one, including at an endpoint.
There is no need to turn them into small nonzero numbers by clipping.
Covariance eigenvalues determine relative many-body singular values,
while the record probability restores their absolute scale. This is an
operator spectrum, distinct from the derivative of normalized conditional
state evolution with respect to its initial state.

The conclusion is that gain/loss feedback admits a controlled Gaussian boundary representation,
and exact covariance propagation remains well defined on each successful
branch. This observation is not a proof of criticality, a determination
of central charge, or an identification of every regulator-stable
one-particle rate with a many-body or tangent Lyapunov exponent.

\bibliographystyle{apsrev4-2}
\bibliography{references}
\end{document}
